\section{Path receipts and gauge}

For a directed path
\[
P=a_1a_2\cdots a_k,
\]
define its ordered receipt by
\[
\alpha(P)=\alpha(a_1)\alpha(a_2)\cdots\alpha(a_k).
\]
Lifting $P$ from $(t(a_1),x)$ ends at
\[
(h(a_k),x\alpha(P)).
\]
The reverse path has receipt $\alpha(P)^{-1}$.

A vertex gauge is a function $g:V(X)\to R$. It changes fiber coordinates by
\[
G_g(v,x)=(v,xg(v)).
\]
Using the convention that the transformed lift is
\[
\widetilde X_{\alpha^g}=G_g\widetilde X_{\alpha}G_g^{-1},
\]
the transformed receipt on a dart $a:u\to v$ is
\[
\alpha^g(a)=g(u)^{-1}\alpha(a)g(v).
\]

\begin{lemma}[Path telescoping]
For every path $P$ from $u$ to $v$,
\[
\alpha^g(P)=g(u)^{-1}\alpha(P)g(v).
\]
\end{lemma}

\begin{proof}
Insert the edge transformation formula into the ordered path product. At
every internal vertex, the adjacent factors $g(w)g(w)^{-1}$ cancel.
\end{proof}

If $P$ is a loop based at $r$, its receipt therefore transforms as
\[
\alpha^g(P)=g(r)^{-1}\alpha(P)g(r).
\]
Gauge preserves the conjugacy class of every based circuit receipt and the
conjugacy class of the subgroup generated by all such receipts.
